\SourceLeaf{231}%
% A544-P0231-D01
ingredi $(n)$, unde patet, quomodo omnes quantitates per praecedentes sint determinandae. Ita invenitur
\[
\begin{array}{lll}
(1)=p&(4)=\frac14(p^4-pp)&(7)=\frac17(p^7-p)\\
(2)=\frac12(pp-p)&(5)=\frac15(p^5-p)&(8)=\frac18(p^8-p^4)\\
(3)=\frac13(p^3-p)&(6)=\frac16(p^6-p^3-pp+p)&\text{etc.}
\end{array}
\]

% A544-P0231-D02
\Article{344--346}
Observatur ex hoc seriei initio, summum terminum expressionis $(n)$ esse $\frac1n p^n$, ad quem, si $n$ est primus, accedit $-\frac1n p$; at si $n$ est compositus, lex minus elucet. Si vero attentius rem consideramus, videmus esse
\[
\begin{array}{ll}
p=(1)&p^5=5(5)+(1)\\
pp=2(2)+(1)&p^6=6(6)+3(3)+2(2)+(1)\\
p^3=3(3)+(1)&p^7=7(7)+(1)\\
p^4=4(4)+2(2)+(1)&p^8=8(8)+4(4)+2(2)+(1)\text{ etc.}
\end{array}
\]
ubi lex progressionis est manifesta; scilicet si omnes numeri $n$ divisores sint $\alpha,\beta,\gamma,\delta$ etc., erit
\[
p^n=\alpha(\alpha)+\beta(\beta)+\gamma(\gamma)+\delta(\delta)+\text{etc.}
\]

% A544-P0231-D03
Huius observationis generalitatem iam demonstrare accingimur.

% A544-P0231-D04
Ostendimus summam omnium talium expressionum $(1^\alpha)(2^\beta)(3^\gamma)\ldots$ si semper $\alpha+2\beta+3\gamma+\ldots=n$, exhaurire omnes functiones $n$ dimensionum adeoque esse $=p^n$. Hinc patet, \textemdash\ \textemdash\ \textemdash. Si
\[
\left(\frac1{1-x}\right)^{(1)}\left(\frac1{1-xx}\right)^{(2)}\left(\frac1{1-x^3}\right)^{(3)}\ldots
\]
evolvatur in seriem $1+Ax+Bx^2\ldots=P$, erit
\[
A=p,\quad B=p^2,\quad C=p^3\text{ etc.}
\]
\[
\frac{x\,\mathrm{d}P}{P\,\mathrm{d}x}=\frac{(1)x}{1-x}+\frac{2(2)x^3}{1-x^2}+\frac{3(3)x^3}{1-x^3}\ldots
\]
\noindent\hspace*{.28\linewidth}\textemdash\ \textemdash\ \textemdash\ \textemdash

% A544-P0231-D05
[hinc substituendo $\frac{px}{1-px}$ pro $\frac{x\,\mathrm{d}P}{P\,\mathrm{d}x}$ et evolvendo singulas fractiones in series infinitas theorematis veritas sponte elucet.]

\SourceLeaf{232}%
% A544-P0232-D01
\Article{347}
Theorema hoc etiam alio modo exprimi potest. Scilicet si numeri $n$ divisores omnes sint $n,1,\delta,\delta',\delta'',\delta'''$ etc., theorema in eo consistit, ut sit
\[
p^n=n(n)+(1)+\delta(\delta)+\delta'(\delta')+\text{etc.}
\]
Iam patet, productum ex $(n)$ functionibus primis, quae sunt $n$ dimensionum, habere $n(n)$ dimensiones et sic de reliquis, quare

% A544-P0232-D02
\emph{Productum ex omnibus functionibus primis dimensionis unius, dimensionum $n,\delta,\delta'$ etc. habebit $p^n$ dimensiones.}

% A544-P0232-D03
Facile nunc est ex hoc theoremate valorem expressionis $(n)$ ipsam deducere, sed brevitatis gratia analysin, quae non est difficilis, supprimimus. Sit itaque $n=a^\alpha b^\beta c^\gamma$ etc., ita ut $a,b,c$ etc. sint numeri primi diversi, eritque
\[
n(n)=p^n-\Sigma p^{\frac na}+\Sigma p^{\frac n{ab}}-\Sigma p^{\frac n{abc}}\text{ etc.}
\]
ubi $\Sigma p^{\frac n{abc\dotstwo}}$ significat complexum omnium expressionum huic $p^{\frac n{abc\dotstwo}}$ similium, si quantitates $a,b,c\dotstwo$ quomodocunque inter se permutentur. Ita pro $n=36$ erit $36(36)=p^{36}-p^{18}-p^{12}+p^6$.

% A544-P0232-D04
Unam adhuc observationem adiicere liceat. Si $n$ est formae $a^\alpha$ et $a$ primus, erit $n(n)=p^n-p^{\frac na}$, quare, quum $(n)$ necessario sit integer, erit quicquid sit $p$,
\[
p^n\equiv p^{\frac na}\;(\mathrm{mod.}\,n)
\]
quare, si $p$ ad $a$ primus erit,
\[
p^{n-\frac na}\equiv1\;(\mathrm{mod.}\,n)
\]
et pro $\alpha=1$
\[
p^{a-1}\equiv1\;(\mathrm{mod.}\,a)
\]
Memorabile est, haec theoremata tam diversis modis erui posse.

% A544-P0232-D05
\Article{348}
\textsc{Problema.} \emph{Data aequatione}
\[
x^m+Ax^{m-1}+Bx^{m-2}+Cx^{m-3}+\text{etc.}+M=0
\]
\emph{cuius radices sunt $x=a$, $x=b$, $x=c$ etc., invenire aequationem, cuius radices sint $x=a^\tau$, $x=b^\tau$, $x=c^\tau$ etc.}

\SourceLeaf{233}%
% A544-P0233-D01
\emph{Solutio prima.} Quaerantur per theorema notum summae radicum aequationis propositae, earum quadratorum, cuborum etc. usque ad potestatem $m\tau^{\mathrm{tam}}$. Hinc igitur habentur etiam summae radicum aequationis quaesitae nec non quadratorum etc. scilicet $\Sigma a^\tau$, $\Sigma a^{2\tau}$ etc., unde per idem theorema coëfficientes determinari possunt.

% A544-P0233-D02
Ad praxin quidem haec solutio est facilior; sed ad institutum nostrum nec non ad ostendendum, coëfficientes aequationis quaesitae fore integros, si aequationis propositae coëfficientes fuerint integri, quae sequitur magis est accomodata.

% A544-P0233-D03
\emph{Solutio secunda.} Sit $\theta$ radix prima aequationis $x^\tau=1$, fiatque productum ex
\[
\begin{aligned}
&x^m+Ax^{m-1}+Bx^{m-2}+\text{etc.}\\
&x^m+A\theta x^{m-1}+B\theta\theta x^{m-2}+\text{etc.}\\
&x^m+A\theta\theta x^{m-1}+B\theta^4x^{m-2}+\text{etc.}\\
&\text{etc.}\\
&x^m+A\theta^{\tau-1}x^{m-1}+B\theta^{2\tau-2}x^{m-2}+\text{etc.}
\end{aligned}
\]
Huius itaque producti radices erunt
\[
\begin{array}{llll}
a,&\theta a,&\theta\theta a&\text{etc.}\\
b,&\theta b,&\theta\theta b&\text{etc.}\\
c,&\theta c,&\theta\theta c&\text{etc.}
\end{array}
\]
i. e. productum aequale erit huic
\[
(x^\tau-a^\tau)(x^\tau-b^\tau)(x^\tau-c^\tau)\ldots
\]
adeoque huius formae
\[
x^{\tau m}+A'x^{\tau(m-1)}+B'x^{\tau(m-2)}+\text{etc.}
\]
Iam si pro $x^\tau$ scribatur $x$, erit
\[
x^m+A'x^{m-1}+B'x^{m-2}+\text{etc.}=(x-a^\tau)(x-b^\tau)(x-c^\tau)\ldots
\]
adeoque
\[
x^m+A'x^{m-1}+B'x^{m-2}+\text{etc.}=0
\]
aequatio quaesita. Quod vero hic $A',B'$ etc. sint non solum rationales sed etiam integri, facile ex theoria aequationis $x^\tau=1$ deducitur.

% A544-P0233-D04
Quoniam hac operatione in sequentibus saepe utemur, per $(P,\rho^\tau)$ 
\SourceLeaf{234}%
% A544-P0234-D01
indicabimus functionem, qua cifrae aequali posita aequatio proveniens habeat radices, quae sunt potestates $\tau^{\mathrm{tae}}$ radicum aequationis $P=0$.

% A544-P0234-D02
Si $P\equiv Q$ secundum modulum quemcunque, erit etiam $(P,\rho^\tau)\equiv(Q,\rho^\tau)$ secundum eundem modulum.

% A544-P0234-D03
\Article{349}
\textsc{Theorema.} \emph{Coëfficiens termini $x^n$ in $(P,\rho^\tau)$ congruus est secundum modulum $\tau$ coëfficiens termini $x^{\tau n}$ in $P^\tau$, siquidem $\tau$ est numerus primus} (quod pro hoc casu est tertia solutio problematis praecedentis).

% A544-P0234-D04
\emph{Demonstr.} Ex capite sexto sequitur, producti
\[
(x^m+Ax^{m-1}+\text{etc.})(x^m+A\theta x^{m-1}+\text{etc.})\ldots
\]
coëfficientem quemcunque habere hanc formam, postquam pro $\theta^\tau$ substituta est unitas,
\[
E+(1+\theta+\theta\theta+\text{etc.}+\theta^{\tau-1})F
\]
Quodsi iam $\theta$ consideretur tamquam radix prima aequationis $x^\tau=1$, totum productum abiit in $E$; si vero ponatur $\theta=1$, totum productum abiit in $P^\tau=E+\tau F$, quare erit coëfficiens termini $x^{n\tau}$ in $P^\tau$ congruus secundum modulum $\tau$ coëfficienti termini $x^{n\tau}$ in $E$, i. e. coëfficienti termini $x^n$ in $(P,\rho^\tau)$.

% A544-P0234-D05
\Article{350}
\textsc{Theorema.} \emph{Si $\tau$ est numerus primus, erit}
\[
(P,\rho^\tau)\equiv P\;(\mathit{mod.}\,\tau)
\]

% A544-P0234-D06
\emph{Demonstr.} Sit coëfficiens termini $x^n$ in $(P,\rho^\tau)=N'$, in $P$ vero eiusdem termini coëfficiens $=N$. Tunc posito
\[
P=x^m+Ax^{m-1}+\text{etc.}+Nx^n+\text{etc.}
\]
erit
\[
P^\tau\equiv x^{m\tau}+A^\tau x^{(m-1)\tau}+\text{etc.}+N^\tau x^{n\tau}+\text{etc.}\;(\mathrm{mod.}\,\tau)
\]
adeoque (§. praec.) $N'\equiv N^\tau\;(\mathrm{mod.}\,\tau)$; quare, quum $N^\tau\equiv N$, erit $N'\equiv N$. Q.\,E.\,D.

% A544-P0234-D07
Hinc etiam patet, esse $(P,\rho^\alpha)\equiv(P,\rho^{\alpha\tau})$ et $(P,\rho^\tau)\equiv(P,\rho^{\tau\tau})$, unde generaliter
\[
(P,\rho^\alpha)\equiv(P,\rho^{\alpha\tau^k})\;(\mathrm{mod.}\,\tau)
\]

\SourceLeaf{235}%
% A544-P0235-D01
\Article{351}
\textsc{Theorema.} \emph{Datur valor numeri $\nu$ minor quam $p^m$, ita ut functio $x^\nu-1$ per functionem propositam $P$ $m$ dimensionum, cuius pars infima indeterminatam $x$ non involvit, secundum modulum $p$ dividi possit.}

% A544-P0235-D02
\emph{Dem.} Dividatur per $P$ series functionum $1,x,xx\ldots$ usque ad $x^{p^m-1}$, simulac dimensionem $m$ superant, et quoniam nulla per $P$ sine residuo dividi poterit, omnia residua ad hanc formam redigi poterunt
\[
Ax^{m-1}+Bx^{m-2}+\dotstwo+N
\]
ita ut omnes coëfficientes sint positivi et $<p$. Sed patet, quum nunquam omnes possint esse $=0$, $p^m-1$ tantummodo functiones dari, quarum alicui singulae aequales esse debent, quare quum usque ad potestatem ipsius $x$, cuius exponens est $p^m-1$, $p^m$ residua habeantur, necessario duo ad minimum eadem esse debent. Prodeat igitur idem residuum, si $x^a$ et $x^{a+\nu}$ per $P$ dividantur, ita ut $a+\nu<p^m$. Quare $x^{a+\nu}-x^a$ per $P$ dividi poterit. Hinc quoniam (hyp.) $x$ adeoque etiam $x^a$ functio est ad $P$ prima, etiam $x^\nu-1$ per $P$ dividi poterit Q.\,E.\,D.

% A544-P0235-D03
\emph{Coroll.} Si $x^\nu-1$ per $P$ dividatur, etiam $x^{k\nu}-1$ per $P$ dividi poterit, denotante $k$ numerum quemcunque integrum.

% A544-P0235-D04
\Article{352}
\textsc{Theorema.} \emph{Manentibus denominationibus ut in §. praec., si $P$ fuerit functio prima et $x^\nu$ infima potestas, quae unitate mulctata per $P$ dividi possit, erit $\nu$ aut $=p^m-1$ aut pars aliquota huius numeri, excepto unico casu, ubi $P=x$.}

% A544-P0235-D05
\emph{Dem.} Quoniam $P$ est functio prima $m$ dimensionum, dabuntur $p^m-1$ functiones diversae pauciorum quam $m$ dimensionum (exclusa scilicet ab omnium numero functione $0$), quae omnes ad $P$ erunt primae. Iam quum $x^\nu$ supponatur esse infima potestas, quae per $P$ divisa unitatem relinquit, palam est, si omnes inferiores potestates ab $1,x,\dotstwo$ usque ad $x^{\nu-1}$ per $P$ dividantur, $\nu$ residua diversa prodire, quae per $A$ generaliter designentur. Iam si haec exhauriant omnia quae sunt possibilia, theorema erit demonstratum; sin vero quaedam nondum sint in eorum numero, sit quodcunque eorum $B$; iam perspicuum est, functionem $Bx^\nu$ per $P$ divisam residuum $B$ dare et generaliter esse $Bx^{\nu+k}\equiv Bx^k(\mathrm{mod.}\,P)$; sed omnes functiones ab $B$ usque ad $Bx^{\nu-1}$ diversae inter se et ab residuis $A$
